\documentclass{beamer}
\usepackage[utf8]{inputenc}

\usetheme{Madrid}
\usecolortheme{default}
\usepackage{amsmath,amssymb,amsfonts,amsthm}
\usepackage{txfonts}
\usepackage{tkz-euclide}
\usepackage{listings}
\usepackage{adjustbox}
\usepackage{array}
\usepackage{tabularx}
\usepackage{gvv}
\usepackage{lmodern}
\usepackage{circuitikz}
\usepackage{tikz}
\usepackage[utf8]{inputenc}

\usetheme{Madrid}
\usecolortheme{default}
\usepackage{amsmath,amssymb,amsfonts,amsthm}
\usepackage{txfonts}
\usepackage{tkz-euclide}
\usepackage{listings}
\usepackage{adjustbox}
\usepackage{array}
\usepackage{tabularx}
\usepackage{gvv}
\usepackage{lmodern}
\usepackage{circuitikz}
\usepackage{tikz}
\usepackage{graphicx}

\setbeamertemplate{page number in head/foot}[totalframenumber]

\usepackage{tcolorbox}
\tcbuselibrary{minted,breakable,xparse,skins}



\definecolor{bg}{gray}{0.95}
\DeclareTCBListing{mintedbox}{O{}m!O{}}{%
  breakable=true,
  listing engine=minted,
  listing only,
  minted language=#2,
  minted style=default,
  minted options={%
    linenos,
    gobble=0,
    breaklines=true,
    breakafter=,,
    fontsize=\small,
    numbersep=8pt,
    #1},
  boxsep=0pt,
  left skip=0pt,
  right skip=0pt,
  left=25pt,
  right=0pt,
  top=3pt,
  bottom=3pt,
  arc=5pt,
  leftrule=0pt,
  rightrule=0pt,
  bottomrule=2pt,

  colback=bg,
  colframe=orange!70,
  enhanced,
  overlay={%
    \begin{tcbclipinterior}
    \fill[orange!20!white] (frame.south west) rectangle ([xshift=20pt]frame.north west);
    \end{tcbclipinterior}},
  #3,
}
\lstset{
    language=C,
    basicstyle=\ttfamily\small,
    keywordstyle=\color{blue},
    stringstyle=\color{orange},
    commentstyle=\color{green!60!black},
    numbers=left,
    numberstyle=\tiny\color{gray},
    breaklines=true,
    showstringspaces=false,
}
%------------------------------------------------------------
%This block of code defines the information to appear in the
%Title page
\title %optional
{2.10.43}
\date{August  2025}
%\subtitle{A short story}

\author % (optional)
{Bhoomika L - EE25BTECH11014}

\begin{document}


\frame{\titlepage}
\begin{frame}{Question}
Find the direction and normal vectors of the line $y = x + 2$
\end{frame}

\begin{frame}{Solution}
\begin{center}
      The equation of the line can be given by
\end{center}
  \begin{align}
      \Vec{x}=\myvec{x\\y}\\
      \Vec{n}^T=\myvec{1&-1}\\
      c=-2
      \end{align}
The equation of line can be expessed as
\begin{align}
     \Vec{n}^T \cdot\Vec{x}=c
\end{align}
where n is the normal vector of the given line.
\end{frame}

\begin{frame}{Solution}
The direction vector of the line is perpendicular to $\Vec{n^T}$
\begin{align}
  Let \Vec{m}=\myvec{1\\m}
\end{align}

We know that,
\begin{align}
    \Vec{n^T}\cdot \Vec{m}=0\\
    \implies \myvec{1&-1}\myvec{1\\m}=0\\
    1-m=0\\
    m=1
\end{align}
\begin{center}
    $\therefore$ The normal vector of the line is $\Vec{n^T}= \myvec{1&-1}$ and the direction vector $\Vec{m}=\myvec{1\\1}$.
\end{center}
\end{frame}

\begin{frame}[fragile]
\frametitle{C Code }
\begin{lstlisting}
    #include <stdio.h>

typedef struct {
    double x;
    double y;
} Vector2D;

// Function to compute direction and normal vectors
void computeVectors(Vector2D p1, Vector2D p2, Vector2D *direction, Vector2D *normal) {
    direction->x = p2.x - p1.x;
    direction->y = p2.y - p1.y;

    normal->x = -(direction->y);
    normal->y = direction->x;
}

\end{lstlisting}
\end{frame}

\begin{frame}[fragile]
\frametitle{Python Code }
\begin{lstlisting}
import ctypes
import matplotlib.pyplot as plt
import numpy as np

# Define the Vector2D struct in Python to match the C struct
class Vector2D(ctypes.Structure):
    _fields_ = [("x", ctypes.c_double),
                ("y", ctypes.c_double)]

lib = ctypes.CDLL("./code6.so")

# Set argument and return types for the function
lib.computeVectors.argtypes = (Vector2D,Vector2D,ctypes.POINTER(Vector2D), ctypes.POINTER(Vector2D))
lib.computeVectors.restype = None

\end{lstlisting}
\end{frame}

\begin{frame}[fragile]
\frametitle{Python Code }
\begin{lstlisting}
# Create input points
p1 = Vector2D(1.0, 2.0)
p2 = Vector2D(4.0, 6.0)

# Prepare output vectors
direction = Vector2D()
normal = Vector2D()

# Call the C function
lib.computeVectors(p1, p2, ctypes.byref(direction), ctypes.byref(normal))

print(f"Direction Vector: ({direction.x:.2f}, {direction.y:.2f})")
print(f"Normal Vector: ({normal.x:.2f}, {normal.y:.2f})")

\end{lstlisting}
\end{frame}

\begin{frame}[fragile]
\frametitle{Python Code }
\begin{lstlisting}
# Line equation: y = x + 2
# Normal vector = (1, -1)
# Direction vector = (1, 1)

# Define range for x
x_vals = np.linspace(-5, 5, 200)
y_vals = x_vals + 2   # original line

# Perpendicular line: slope = -1 (since slope of y=x+2 is 1)
# Pass it through the y-intercept (0,2)
x_perp = np.linspace(-5, 5, 200)
y_perp = -x_perp + 2  # equation of perpendicular

\end{lstlisting}
\end{frame}

\begin{frame}[fragile]
\frametitle{Python Code }
\begin{lstlisting}
# Plot
plt.figure(figsize=(6,6))
plt.plot(x_vals, y_vals, 'r-')
plt.plot(x_perp, y_perp, 'b-')

# Mark intercept
plt.scatter(0, 2, color='black', zorder=5)
plt.text(0.2, 2.1, '(0,2)', fontsize=9)

# Show direction and normal vectors at point (0,2)
plt.arrow(0, 2, 1, 1, head_width=0.2, head_length=0.3, 
          fc='red', ec='red')
\end{lstlisting}
\end{frame}

\begin{frame}[fragile]
\frametitle{Python Code }
\begin{lstlisting}
plt.arrow(0, 2, 1, -1, head_width=0.2, head_length=0.3, 
          fc='blue', ec='blue')

# --- Label the original line with its equation directly on the line ---
plt.text(1, 1+2+0.2, r"$y = x + 2$", color='red', fontsize=11)

# --- Add text labels in the top-right corner ---
plt.text(3.5, 7, "Red – Original Line (y = x + 2)", color='red', fontsize=10)
plt.text(3.5, 6.3, "Blue – Perpendicular Line (y = -x + 2)", color='blue', fontsize=10)
\end{lstlisting}
\end{frame}

\begin{frame}[fragile]
\frametitle{Python Code }
\begin{lstlisting}
# Axes and grid
plt.axhline(0, color='gray', linewidth=0.5)
plt.axvline(0, color='gray', linewidth=0.5)
plt.gca().set_aspect('equal', adjustable='box')
plt.grid(True)

# Save as fig_6.png
plt.savefig("fig_6.png", dpi=300)
plt.show()
\end{lstlisting}
\end{frame}

\begin{frame}{Plot}
    \centering
    \includegraphics[width=\columnwidth, height=0.8\textheight, keepaspectratio]{../Figs/fig_6.png}     
\end{frame}
\end{document}
